\documentclass{article}
\usepackage{amsmath}
\usepackage{amsfonts}
\usepackage{esint}
\usepackage[UTF8]{ctex}
\usepackage{multirow}
\usepackage{longtable}
\usepackage[left=1cm,right=1cm]{geometry}
\renewcommand{\arraystretch}{1.8}

\title{\huge{Terms in General Physics}}
\author{Huang Zixuan 20250065}
\date{\today}

\begin{document}

  \maketitle

  \section{牛顿力学}
\begin{longtable}{|c|c|c|c|}
  \caption{Terms in Newtonian Mechanics (牛顿力学术语表)} \label{Tab.1} \\
  \hline
  Part (模块) & Terms (英语术语) & Translations (中文翻译) & Expressions (物理记号) \\ 
  \hline
  \endfirsthead

  \hline
  Part (模块) & Terms (英语术语) & Translations (中文翻译) & Expressions (物理记号) \\ 
  \hline
  \endhead

  \hline
  \multicolumn{4}{r}{续表} \\
  \endfoot

  \hline
  \endlastfoot

  \multirow{8}{*}{质点运动学}
  & position vector & 位矢 & $\boldsymbol{r}$ \\
  \cline{2-4}
  & displacement & 位移 & $\Delta\boldsymbol{r}$ \\
  \cline{2-4}
  & velocity & 速度 & $\boldsymbol{v}=\frac{d\boldsymbol{r}}{dt}$ \\
  \cline{2-4}
  & speed & 速率 & $v=|\boldsymbol{v}|$ \\
  \cline{2-4}
  & acceleration & 加速度 & $\boldsymbol{a}=\frac{d\boldsymbol{v}}{dt}$ \\
  \cline{2-4}
  & tangential acceleration & 切向加速度 & $a_{\tau}=\frac{dv}{dt}$ \\
  \cline{2-4}
  & normal (centripetal) acceleration & 法向（向心）加速度 & $a_n=\frac{v^2}{\rho}$ \\
  \cline{2-4}
  & angular position & 角位置 & $\theta$ \\
  \hline
  \multirow{4}{*}{质点运动学}
  & angular velocity & 角速度 & $\boldsymbol{\omega}=\frac{d\boldsymbol{\theta}}{dt}$ \\
  \cline{2-4}
  & angular acceleration & 角加速度 & $\boldsymbol{\alpha}=\frac{d\boldsymbol{\omega}}{dt}$ \\
  \cline{2-4}
  & period & 周期 & $T$ \\
  \cline{2-4}
  & frequency & 频率 & $\nu \quad f$ \\
  \hline

  \multirow{10}{*}{质点动力学}
  & mass & 质量 & $m$ \\
  \cline{2-4}
  & force & 力 & $\boldsymbol{F}$ \\
  \cline{2-4}
  & Newton's first law & 牛顿第一定律 & / \\
  \cline{2-4}
  & Newton's second law & 牛顿第二定律 & $\boldsymbol{F}=\frac{d\boldsymbol{p}}{dt}=m\boldsymbol{a}$ \\
  \cline{2-4}
  & Newton's third law & 牛顿第三定律 & $\boldsymbol{F}_{12}=-\boldsymbol{F}_{21}$ \\
  \cline{2-4}
  & inertial reference frame & 惯性参考系 & / \\
  \cline{2-4}
  & non-inertial reference frame & 非惯性参考系 & / \\
  \cline{2-4}
  & inertial force & 惯性力 & $\boldsymbol{F}_{\text{in}}=-m\boldsymbol{a}_{\text{ref}}$ \\
  \cline{2-4}
  & centripetal force & 向心力 & $\boldsymbol{F}_n=m\frac{v^2}{\rho}\hat{n}$ \\
  \cline{2-4}
  & centrifugal force & 离心力 & / \\
  \hline

  \multirow{3}{*}{功与能}
  & work & 功 & $W=\int_{A}^{B}\boldsymbol{F}\cdot d\boldsymbol{r}$ \\
  \cline{2-4}
  & kinetic energy & 动能 & $E_k=\frac{1}{2}mv^2$ \\
  \cline{2-4}
  & work-energy theorem & 动能定理 & $W=\Delta E_k$ \\
  \hline
  \multirow{9}{*}{功与能}
  & potential energy & 势能 & $E_p$ \\
  \cline{2-4}
  & gravitational potential energy & 重力势能 & $E_p=mgh$ \\
  \cline{2-4}
  & elastic potential energy & 弹性势能 & $E_p=\frac{1}{2}kx^2$ \\
  \cline{2-4}
  & conservative force & 保守力 & $\boldsymbol{F}=-\nabla E_p$ \\
  \cline{2-4}
  & non-conservative force & 非保守力 & / \\
  \cline{2-4}
  & mechanical energy & 机械能 & $E=E_k+E_p$ \\
  \cline{2-4}
  & conservation of mechanical energy & 机械能守恒 & $\Delta E=0$ \\
  \cline{2-4}
  & power & 功率 & $P=\frac{dW}{dt}=\boldsymbol{F}\cdot\boldsymbol{v}$ \\
  \cline{2-4}
  & efficiency & 效率 & $\eta$ \\
  \hline

  \multirow{8}{*}{动量与角动量}
  & linear momentum & 动量 & $\boldsymbol{p}=m\boldsymbol{v}$ \\
  \cline{2-4}
  & impulse & 冲量 & $\boldsymbol{I}=\int_{t_1}^{t_2}\boldsymbol{F}\,dt$ \\
  \cline{2-4}
  & impulse-momentum theorem & 动量定理 & $\boldsymbol{I}=\Delta\boldsymbol{p}$ \\
  \cline{2-4}
  & conservation of linear momentum & 动量守恒定律 & $\Sigma\boldsymbol{p}_i=\text{const}$ \\
  \cline{2-4}
  & center of mass & 质心 & $\boldsymbol{r}_C=\frac{\Sigma m_i\boldsymbol{r}_i}{\Sigma m_i}$ \\
  \cline{2-4}
  & angular momentum & 角动量 & $\boldsymbol{L}=\boldsymbol{r}\times\boldsymbol{p}$ \\
  \cline{2-4}
  & torque & 力矩 & $\boldsymbol{M}=\boldsymbol{r}\times\boldsymbol{F}$ \\
  \cline{2-4}
  & conservation of angular momentum & 角动量守恒 & $\boldsymbol{L}=\text{const}$ \\
  \hline

  \multirow{10}{*}{刚体力学}
  & rigid body & 刚体 & / \\
  \cline{2-4}
  & moment of inertia & 转动惯量 & $I=\Sigma m_i r_i^2$ \\
  \cline{2-4}
  & parallel-axis theorem & 平行轴定理 & $I=I_C+md^2$ \\
  \cline{2-4}
  & perpendicular-axis theorem & 垂直轴定理 & $I_z=I_x+I_y$ \\
  \cline{2-4}
  & rotational kinetic energy & 转动动能 & $E_{kr}=\frac{1}{2}I\omega^2$ \\
  \cline{2-4}
  & torque (rotational dynamics) & 力矩（转动动力学） & $\boldsymbol{M}=I\boldsymbol{\alpha}$ \\
  \cline{2-4}
  & angular momentum (rigid body) & 角动量（刚体） & $\boldsymbol{L}=I\boldsymbol{\omega}$ \\
  \cline{2-4}
  & rolling motion & 滚动（纯滚动） & / \\
  \cline{2-4}
  & gyroscope & 陀螺仪 & / \\
  \cline{2-4}
  & precession & 进动 & / \\
  \hline
\end{longtable}

\section{热学}
\begin{longtable}{|c|c|c|c|}
  \caption{Terms in Thermal Physics (热学术语表)} \label{Tab.2} \\
  \hline
  Part (模块) & Terms (英语术语) & Translations (中文翻译) & Expressions (物理记号) \\ 
  \hline
  \endfirsthead

  \hline
  Part (模块) & Terms (英语术语) & Translations (中文翻译) & Expressions (物理记号) \\ 
  \hline
  \endhead

  \hline
  \multicolumn{4}{r}{续表} \\
  \endfoot

  \hline
  \endlastfoot

  \multirow{2}{*}{热力学}
  & temperature & 温度 & $T$ \\
  \cline{2-4}
  & thermodynamic equilibrium & 热力学平衡 & / \\
  \hline
  \multirow{8}{*}{热力学}
  & zeroth law of thermodynamics & 热力学第零定律 & / \\
  \cline{2-4}
  & equation of state & 状态方程 & $f(p,V,T)=0$ \\
  \cline{2-4}
  & ideal gas equation of state & 理想气体状态方程 & $pV=nRT$ \\
  \cline{2-4}
  & mole & 摩尔 & $\nu$ \\
  \cline{2-4}
  & Avogadro's number & 阿伏伽德罗常量 & $N_A$ \\
  \cline{2-4}
  & Boltzmann's constant & 玻尔兹曼常数 & $k_B$ \\
  \cline{2-4}
  & gas constant & 气体常量 & $R=8.314\,\text{J/(mol·K)}$ \\
  \cline{2-4}
  & degrees of freedom & 自由度 & $i$ \\
  \hline

  \multirow{9}{*}{热力学定律}
  & work (thermodynamics) & 功（热力学） & $W=\int_{V_1}^{V_2}p\,dV$ \\
  \cline{2-4}
  & heat & 热量 & $Q$ \\
  \cline{2-4}
  & internal energy & 内能 & $U$ \\
  \cline{2-4}
  & first law of thermodynamics & 热力学第一定律 & $\Delta U=Q+W$（或 $\Delta U=Q-W$） \\
  \cline{2-4}
  & molar heat capacity & 摩尔热容 & $C_m$ \\
  \cline{2-4}
  & heat capacity at constant pressure & 定压热容 & $C_p$ \\
  \cline{2-4}
  & heat capacity at constant volume & 定体热容 & $C_V$ \\
  \cline{2-4}
  & adiabatic process & 绝热过程 & $pV^\gamma=\text{const}$ \\
  \cline{2-4}
  & isothermal process & 等温过程 & $pV=\text{const}$ \\
  \hline
  \multirow{3}{*}{热力学定律}
  & isobaric process & 等压过程 & $V/T=\text{const}$ \\
  \cline{2-4}
  & isochoric process & 等体过程 & $p/T=\text{const}$ \\
  \cline{2-4}
  & cyclic process & 循环过程 & $\Delta U=0$ \\
  \hline

  \multirow{12}{*}{热力学第二定律}
  & second law of thermodynamics & 热力学第二定律 & / \\
  \cline{2-4}
  & Kelvin-Planck statement & 开尔文-普朗克表述 & / \\
  \cline{2-4}
  & Clausius statement & 克劳修斯表述 & / \\
  \cline{2-4}
  & heat engine & 热机 & / \\
  \cline{2-4}
  & refrigerator & 制冷机 & / \\
  \cline{2-4}
  & efficiency (heat engine) & 热机效率 & $\eta=1-\frac{Q_2}{Q_1}$ \\
  \cline{2-4}
  & coefficient of performance & 制冷系数 & $\varepsilon=\frac{Q_2}{W}$ \\
  \cline{2-4}
  & Carnot cycle & 卡诺循环 & / \\
  \cline{2-4}
  & Carnot theorem & 卡诺定理 & / \\
  \cline{2-4}
  & entropy & 熵 & $S$ \\
  \cline{2-4}
  & Clausius inequality & 克劳修斯不等式 & $\oint\frac{\delta Q}{T}\le 0$ \\
  \cline{2-4}
  & entropy increase principle & 熵增加原理 & $\Delta S\ge 0$ \\
  \hline

  \multirow{2}{*}{气体动理论}
  & kinetic theory of gases & 气体动理论 & / \\
  \cline{2-4}
  & mean free path & 平均自由程 & $\bar{\lambda}$ \\
  \hline
  \multirow{7}{*}{气体动理论}
  & mean collision frequency & 平均碰撞频率 & $\bar{Z}$ \\
  \cline{2-4}
  & Maxwell distribution & 麦克斯韦分布 & / \\
  \cline{2-4}
  & most probable speed & 最概然速率 & $v_p=\sqrt{\frac{2kT}{m}}$ \\
  \cline{2-4}
  & mean speed & 平均速率 & $\bar{v}=\sqrt{\frac{8kT}{\pi m}}$ \\
  \cline{2-4}
  & root-mean-square speed & 方均根速率 & $v_{\text{rms}}=\sqrt{\frac{3kT}{m}}$ \\
  \cline{2-4}
  & pressure & 压强 & $p=\frac{1}{3}nm\bar{v}^2$ \\
  \cline{2-4}
  & equipartition theorem & 能量均分定理 & / \\
  \hline
\end{longtable}

\end{document}